\SourceLeaf{266}
\fancyhead[L]{\small INVESTIGATIONS INTO TOPICS IN HIGHER GEODESY}
\renewcommand{\Article}[1]{\par\medskip\Needspace{4\baselineskip}\pdfbookmark[1]{Art. #1}{geodaesie1-art-#1}\centerline{\textbf{#1.}}\nobreak\smallskip}
\begin{center}
INVESTIGATIONS\par\bigskip
INTO\par\bigskip
{\large TOPICS IN HIGHER GEODESY.}\par\medskip
\rule{.65in}{.3pt}
\end{center}
\noindent\hspace*{\parindent}The trigonometric surveys that extended over the whole Kingdom of Hanover, carried out partly by me and partly under my direction, followed methods that differ from those usually employed. This concerns the way the measurements were made and, still more, their subsequent mathematical treatment and the way the results were obtained from them. I had originally intended, once the measurements were entirely finished, to set them out in a separate work together with every procedure I had used. For reasons I need not go into here, I have so far been unable to carry out that plan. I am therefore choosing to publish what is distinctive about my theoretical work in a series of memoirs. I am all the more willing to take this course because it leaves me free to develop in detail a number of investigations that are of interest in their own right and closely related to the others, even where I made no direct use of them in my surveys. This applies especially to most of the present first memoir.
